\chapter{Decisions with Incomplete State Information}
\label{sec:p17-partial-observation-decisions}

The distribution of an unobserved state, conditional on the information
actually available to the controller, can itself serve as the state
of a decision problem. \AA str\"om develops this construction for
finite hidden states and observations, continuous controls in a compact
set, and a finite horizon \cite[pp.~174--205]{Astrom1965IncompleteInformation}.
Its new state space is an entire probability simplex.
The original five theorems, belief-process lemma and two four-stage
examples are developed below with the observation timing and conditional
expectations explicit.

\section{The controlled process and the information available}

\begin{hbdefinition}{def:p17-astrom-model}{States, controls, observations and timing}
Let the hidden state set be \(X=\{1,\ldots,n\}\), the observation
set \(Y=\{1,\ldots,m\}\), and let \(U\) be a nonempty compact
subset of \(\mathbb R^d\). The horizon \(N\) is a positive integer,
and time runs from \(1\) to \(N\).
For \(t<N\), a control \(u\in U\) has transition matrix
\(P_t(u)=(p_{si,t}(u))\), with nonnegative entries and row
sums one. Every entry is continuous in \(u\).
The observation matrix \(Q=(q_{ij})\) has nonnegative entries
and row sums one. The cost \(g_t(i,u)\) is continuous in \(u\)
for each \(i,t\). Finite state and time sets and compactness
therefore give a finite bound on all these costs.

At stage \(t\), the controller has observed \(Y_1,\ldots,Y_t\)
and knows its earlier controls \(U_1,\ldots,U_{t-1}\).
Write this history as \(H_t\) and its sigma algebra as
\(\mathcal I_t\). It chooses \(U_t\), incurs cost
\(g_t(X_t,U_t)\), then the hidden state moves according to
\(P_t(U_t)\) and observation \(Y_{t+1}\) is produced.
The precise conditional laws are
\[
 \begin{split}
 P(X_{t+1}=i\mid X_t=s,H_t,U_t=u)&=p_{si,t}(u),\\
 P(Y_{t+1}=j\mid X_{t+1}=i,X_t,H_t,U_t)&=q_{ij}.
 \end{split}
\]
An action kernel depends only on \(H_t\); conditional on that
history its private randomization supplies no further information
about the hidden state. Costs are not an extra observation:
if a realized cost is observed and reveals state information,
it must be included in the observation record.

An initial law \(\mu\) for \(X_1\), before \(Y_1\), is input,
and \(P(Y_1=j\mid X_1=i)=q_{ij}\).
If initialization is instead given at \(X_0\) by the source's
\(p^0\), specify any preceding transition and its fixed
initial control, then take \(\mu=p^0P_0(u_0)\).
This makes the first observation's law definite.
The control chosen after \(Y_t\) governs the transition
to \(X_{t+1}\); it cannot change the state already observed.

The source's admissible policies are deterministic maps
\(u_t=c_t(Y_1,\ldots,Y_t)\). Its histories containing past
controls are equivalent, since those controls are recovered
recursively from the observation history and the known policy.
We also allow Borel randomized history policies and prove
optimality against this larger class.
The objective is
\[
                     J^\pi=E^\pi\sum_{t=1}^N g_t(X_t,U_t).
\]
Maximization of rewards is obtained by replacing \(g\) by \(-g\).
\end{hbdefinition}

\begin{hbdefinition}{def:p17-astrom-belief}{The full continuous information state}
Let
\[
 \Delta_n=\{b\in[0,1]^n:\textstyle\sum_i b_i=1\}.
\]
The belief at stage \(t\) is the vector
\(B_t(i)=P(X_t=i\mid\mathcal I_t)\).
For \(b\in\Delta_n\), define
\[
 a_i(b,u,t)=\sum_s b_s p_{si,t}(u),\qquad
 z_i^j(b,u,t)=q_{ij}a_i(b,u,t),\qquad
 \lambda_j(b,u,t)=\sum_i z_i^j(b,u,t).
\]
Fix a default \(b^\dagger\in\Delta_n\), and put
\[
 F_j(b,u,t)=
 \begin{cases}
 z^j(b,u,t)/\lambda_j(b,u,t),&\lambda_j(b,u,t)>0,\\
 b^\dagger,&\lambda_j(b,u,t)=0.
 \end{cases}
\]
The source writes \(w\) for the belief and
\(\|z^j\|_1\) for \(\lambda_j\).
For the initial observation put
\(\lambda_j^0=\sum_i\mu_iq_{ij}\), and let
\(F_j^0(i)=\mu_iq_{ij}/\lambda_j^0\) when \(\lambda_j^0>0\);
use \(b^\dagger\) otherwise.
The initial belief law is
\(\sum_j\lambda_j^0\delta_{F_j^0}\).
\end{hbdefinition}

\begin{hbtheorem}{thm:p17-astrom-filter}{proved=annex:p17.astrom.filter}{The recursive posterior, including zero-likelihood observations}
The specified laws and any admissible policy determine a unique
joint law. Versions of the conditional beliefs satisfy
\[
 B_1=F_{Y_1}^0,\qquad B_{t+1}=F_{Y_{t+1}}(B_t,U_t,t)
                         \quad\text{almost surely}.
\]
Given \(H_t,U_t\), the next observation has probabilities
\(\lambda_j(B_t,U_t,t)\).
The maps \(F_j,\lambda_j\) are Borel, and the belief transition
kernel is
\[
 {\cal R}_t(C\mid b,u)=
              \sum_{j=1}^m\lambda_j(b,u,t)\mathbf1_C(F_j(b,u,t)).
\]
It is a probability kernel on \(\Delta_n\), supported on at most
\(m\) points. Zero-likelihood branches contribute zero regardless
of their default belief. They are not assigned a posterior by division.
\end{hbtheorem}

\begin{hbtheorem}{thm:p17-astrom-information-state}{proved=annex:p17.astrom.information.state}{The source's belief lemma and the policy bridge}
Under a Borel belief policy \(U_t=f_t(B_t)\), the process \(B_t\)
is a time-dependent Markov process with transition
\({\cal R}_t(\cdot\mid b,f_t(b))\).
For an arbitrary history policy it is a controlled Markov
information state, with this transition conditional on
\(H_t,U_t\), and expected immediate cost
\[
                        \ell_t(b,u)=\sum_i b_i g_t(i,u).
\]
Every belief-history policy can be implemented from the original
observations by the filter, with identical belief/action laws and
expected costs. Conversely, for any original randomized history
policy and a fixed initial law, there is a randomized Markov
policy on beliefs with identical stagewise \((B_t,U_t)\) laws
and total expected cost.
The Markov projection can depend on that initial law.
The finite-horizon optimum below is attained by one deterministic
belief policy for every starting belief, a separate stronger assertion.
\end{hbtheorem}

\section{Existence, the Bellman equation and verification}

\begin{hbdefinition}{def:p17-astrom-recursion}{The original finite-horizon recursion}
Set \(V_{N+1}=0\) and \(\mathcal H_N(b,u)=\ell_N(b,u)\).
Recursively define, on the entire simplex,
\[
 \begin{split}
 \mathcal H_t(b,u)&=\ell_t(b,u)+
       \sum_{j:\lambda_j(b,u,t)>0}
               \lambda_j(b,u,t)V_{t+1}(F_j(b,u,t)),\quad t<N,\\
 V_t(b)&=\min_{u\in U}\mathcal H_t(b,u),\qquad 1\leq t\leq N .
 \end{split}
\]
At \(t=N\) the continuation term is zero, so no terminal
transition matrix is needed; for bookkeeping take
\(B_{N+1}=b^\dagger\).
Equivalently use the sum over all
\(j\) with the explicit default, because its zero weights vanish.
The initial minimum cost is
\[
                       \sum_j\lambda_j^0 V_1(F_j^0).
\]
\end{hbdefinition}

\begin{hblemma}{thm:p17-astrom-continuous-selection}{proved=annex:p17.astrom.continuous.selection}{Continuous weighted posteriors and an attained Borel minimum}
If \(v\) is continuous on \(\Delta_n\), the function
\[
             (b,u)\longmapsto
                 \sum_j\lambda_j(b,u,t)v(F_j(b,u,t))
\]
is continuous, including at vanishing observation probabilities.
For every continuous \(H:\Delta_n\times U\to\mathbb R\),
the function \(b\mapsto\min_uH(b,u)\) is continuous, and
there is a Borel minimizing map \(f:\Delta_n\to U\).
The proof gives a nested compact-cube selector and does not
assume a continuous minimizing map exists.
\end{hblemma}

\begin{hbtheorem}{thm:p17-astrom-finite-optimum}{proved=annex:p17.astrom.finite.optimum}{Complete finite-horizon construction: Theorems 1--3}
The recursion has a unique sequence of bounded continuous solutions.
Every minimum is attained by a Borel selector \(f_t(b)\).
The filtered policy \(U_t=f_t(B_t)\) minimizes total expected
cost against every randomized observation-history policy.
It is optimal from every starting stage and every belief,
including beliefs not reachable from the initially selected \(\mu\).
The original partial-observation problem and its fully observed
belief problem have the same minimum and both have an optimal policy.
This establishes the source's necessary recursion, sufficient
verification theorem and information-state equivalence with
their existence premises supplied.
\end{hbtheorem}

\begin{hbtheorem}{thm:p17-astrom-verification}{proved=annex:p17.astrom.verification}{Exact verification and the scope of necessity}
Define the nonnegative Bellman residual
\(\delta_t(b,u)=\mathcal H_t(b,u)-V_t(b)\).
For every admissible policy and its specified initial law,
\[
 J^\pi-E[V_1(B_1)]
                  =\sum_{t=1}^N E^\pi[\delta_t(B_t,U_t)].
\]
Thus a policy is optimal for that initial law exactly when
\(\delta_t(B_t,U_t)=0\) almost surely for every \(t\).
There is no necessity condition on a history which the policy
reaches with probability zero.
For any policy, its future cost at a history is a conditional
expectation of future costs, rather than the sum of random
future posterior costs without conditioning.
If the chosen actions have residuals at most \(\eta_t\)
on all histories, their total loss is at most \(\sum_t\eta_t\).
\end{hbtheorem}

\section{Perfect information, no further observations and delays}

\begin{hbdefinition}{def:p17-astrom-comparison-values}{The two source comparison problems}
Let \(S_{N+1}(i)=0\), and define
\[
 S_t(i)=\min_{u\in U}
          [g_t(i,u)+\sum_k p_{ik,t}(u)S_{t+1}(k)],
 \qquad L_t(b)=\sum_i b_iS_t(i).
\]
This problem reveals the actual state before every control.
Let \(O_{N+1}=0\), and define
\[
 O_t(b)=\min_{u\in U}
                  [\ell_t(b,u)+O_{t+1}(bP_t(u))].
\]
This problem starts with belief \(b\) and then receives no
additional observations; its optimum can be attained by a
deterministic control schedule.
At stage \(N\) both expressions omit continuation.
All products \(bP_t(u_t)\cdots P_{k-1}(u_{k-1})\) are ordered
in time; different controls or dates do not produce a power
of one common matrix.
\end{hbdefinition}

\begin{hbtheorem}{thm:p17-astrom-information-bounds}{proved=annex:p17.astrom.information.bounds}{Theorems 4--5 and the value of the declared information}
For every stage and belief,
\[
                          L_t(b)\leq V_t(b)\leq O_t(b).
\]
If all observation rows are the same distribution \((r_j)\),
then no new state information is obtained beyond known controls,
and the partial-observation optimum equals \(O_t\).
The source's constant \(q_{ij}\) case is \(r_j=1/m\).
If \(Q\) is the identity, every positive-probability posterior
after an observation is a state point mass, and the source
complete-information recursion is exactly \(S_t\).

For the pre-observation initial law \(\mu\), one has
\[
 \sum_i\mu_i S_1(i)
 \ \leq\ \sum_j\lambda_j^0V_1(F_j^0)
 \ \leq\ \sum_j\lambda_j^0O_1(F_j^0)
 \ \leq\ O_1(\mu).
\]
The third expression allows the first observation and then
commits to a schedule; the last allows no observation at all.
Relative to the last baseline, the value of the specified
observation stream is
\(O_1(\mu)-\sum_j\lambda_j^0V_1(F_j^0)\).
The value of perfect state information is
\(O_1(\mu)-\sum_i\mu_iS_1(i)\).
Both are nonnegative, and the former does not exceed the latter.
After a fixed current belief \(b\), the analogous values are
\(O_t(b)-V_t(b)\) and \(O_t(b)-L_t(b)\).
\end{hbtheorem}

\begin{hbtheorem}{thm:p17-astrom-delays}{proved=annex:p17.astrom.delays}{The source's delayed-information variation}
Suppose the controller at date \(t\) has only observations
through a specified integer date \(0\leq d_t\leq t\), where \(d_t\) is
nondecreasing, and knows all its earlier controls.
The histories with this information define a finite-horizon
decision problem with an attained deterministic optimum.
It can be solved on the posterior of the finite hidden path
\((X_1,\ldots,X_t)\) by updating with newly delivered observations.
This accounts for observations of past states whose laws cannot
be recovered from the current-state marginal alone.
More delivered observations cannot increase minimum cost.
Taking \(d_t=t\) gives the original problem; taking \(d_t=0\)
gives an open-loop schedule.
\end{hbtheorem}

\section{The two original four-stage calculations}

\begin{hbdefinition}{def:p17-astrom-example-one}{A controlled two-state switch}
For the source's Example 1, let
\[
 P(u)=\begin{pmatrix}1-u&u\\u&1-u\end{pmatrix},
 \quad U=[0,1],\qquad
 Q=\begin{pmatrix}q&1-q\\1-q&q\end{pmatrix},
 \quad 0\leq q\leq1.
\]
There are four cost stages, zero cost at stages \(1,2,3\),
and terminal cost \(\mathbf1_{\{X_4=2\}}\).
Write \(b=(x,1-x)\), \(e=\min(q,1-q)\),
\(\alpha=2e-e^2\), \(\zeta=e^2\), and
\(a=(1-u)x+u(1-x)\).
The source's symbols satisfy
\(w=(x,1-x)\), \(q_0=e\), \(\alpha_0=\alpha\), and \(\beta_0=\zeta\).
The next-observation probabilities and posteriors are
\[
 \begin{array}{c|c|c}
 j&\lambda_j&F_j(1)\\ \hline
 1&qa+(1-q)(1-a)&qa/\lambda_1\\
 2&(1-q)a+q(1-a)&(1-q)a/\lambda_2
 \end{array}
\]
with the declared default on zero denominators and
\(F_j(2)=1-F_j(1)\).
\end{hbdefinition}

\begin{hbtheorem}{thm:p17-astrom-example-one}{proved=annex:p17.astrom.example.one}{The complete four-stage switch solution}
On the entire interval \(0\leq x\leq1\),
\[
 \begin{split}
 V_4(x)&=1-x,\\
 V_3(x)&=\min(x,1-x),\\
 V_2(x)&=\min(x,1-x,e),\\
 V_1(x)&=\min\{x,1-x,\alpha x+\zeta(1-x),
                              \alpha(1-x)+\zeta x\}.
 \end{split}
\]
At stage \(3\), choose \(u=1\) if \(x<1/2\), \(u=0\)
if \(x>1/2\), and any \(u\) at equality.
At stages \(1,2\), either endpoint control \(u=0\) or \(u=1\)
is optimal. At stage \(2\), if \(e\leq x\leq1-e\),
every control in \([0,1]\) is optimal.
The proof supplies all four branch sums for \(V_2\), all
nine for \(V_1\), their equality boundaries and the reduction
to the displayed lower envelopes.

For \(q=4/5\) and \(x=1/10\), the four values
\((V_4,V_3,V_2,V_1)\) are
\((9/10,1/10,1/10,9/125)\).
Perfect state information gives \(L_4=1-x\) and
\(L_1=L_2=L_3=0\).
No further observations give \(O_4=1-x\) and
\(O_1=O_2=O_3=\min(x,1-x)\).
\end{hbtheorem}

\begin{hbdefinition}{def:p17-astrom-example-two}{The source's four-action production model}
In Example 2 the criterion is expected profit maximization.
There are four stages and the following exact rational data:
\[
 \begin{array}{c|cc|cc|rr}
 u&p_{11}&p_{12}&p_{21}&p_{22}&g(1,u)&g(2,u)\\ \hline
 1&.5&.5&.4&.6&20&-5\\
 2&.5&.5&.7&.3&17&-8\\
 3&.8&.2&.4&.6&10&-15\\
 4&.8&.2&.7&.3&7&-18
 \end{array}
 \qquad
 Q=\begin{pmatrix}.8&.2\\.3&.7\end{pmatrix}.
\]
The actions are respectively neither advertising nor research,
research alone, advertising alone, and both.
The state indicates successful or unsuccessful production,
and the observation is an imperfect indication of that state.
These are the source's idealized model values, not measured
business outcomes. Its payoff matrix differs from the Howard
example from which it takes the transitions.

For \(x=b_1\), the predicted state-one probabilities are
\[
       a_1=.4+.1x,\quad a_2=.7-.2x,\quad
       a_3=.4+.4x,\quad a_4=.7+.1x.
\]
The next observation has probabilities
\(\lambda_1=.3+.5a_u\), \(\lambda_2=.7-.5a_u\),
with posteriors \(F_1(1)=.8a_u/\lambda_1\) and
\(F_2(1)=.2a_u/\lambda_2\).
Both denominators are positive throughout this model.
\end{hbdefinition}

\begin{hbtheorem}{thm:p17-astrom-example-two}{proved=annex:p17.astrom.example.two}{Exact incomplete-information profits and policy thresholds}
Writing \(V_t(x)\) now for maximum expected future profit,
the original four-stage problem has the following exact values
\[
 \begin{split}
 V_4(x)&=25x-5,\\
 V_3(x)&=\max(27.5x,\ 20x+4.5),\\
 V_2(x)&=\max(27.375x+7.650,\ 20.250x+11.775),\\
 V_1(x)&=\max(27.38875x+15.09250,\
                              20.22250x+19.25875).
 \end{split}
\]
The displayed terminating decimals are exact rationals.
At stage \(4\), action \(1\) is optimal everywhere.
At stages \(3,2,1\), choose action \(2\) below the respective
thresholds
\[
                 \frac35,\qquad\frac{11}{19},
                                      \qquad\frac{1111}{1911},
\]
and action \(1\) above them; either is optimal at equality.
The remaining actions are strictly inferior to one of these
throughout the belief interval. The proof includes every
action's intermediate affine expression.
\end{hbtheorem}

\begin{hbtheorem}{thm:p17-astrom-example-two-bounds}{proved=annex:p17.astrom.example.two.bounds}{The original perfect-information and open-loop comparisons}
For the same profit model, let \(L_t\) now denote the
perfect-information upper value and \(O_t\) the no-further-observation
lower value. Their exact formulas are
\[
 \begin{array}{c|c|c}
 t&L_t(x)&O_t(x)\\ \hline
 4&25x-5&25x-5\\
 3&23x+4.5&\max(27.5x,\ 20x+4.5)\\
 2&23.4x+12.6&\max(27x+7.5,\ 19.5x+11.25)\\
 1&23.32x+20.98&\max(26.95x+14.05,\ 19.6x+18.4).
 \end{array}
\]
The open-loop thresholds at stages \(3,2,1\) are
\(3/5,1/2,29/49\), respectively.
For every \(t,x\), \(O_t(x)\leq V_t(x)\leq L_t(x)\).
At current belief \(x=1/2\), the stage-one values are
\[
                    O_1=28.2,\qquad V_1=29.37,\qquad L_1=32.64.
\]
Thus the specified future measurements add \(1.17\) expected
profit over the schedule baseline; perfect future state
information adds \(4.44\) with that same current belief.
\end{hbtheorem}

\section{Boundaries of the construction}

\begin{hbtheorem}{thm:p17-astrom-information-failures}{proved=annex:p17.astrom.information.failures}{Why the observation and state assumptions matter}
The following explicit finite models show the scope of the method.
\begin{enumerate}
\item An observation with likelihood zero permits any default
      posterior on that event, but permits no likelihood ratio
      with zero denominator.
\item One-time conditional probabilities \(P(Y_t=j\mid X_t=i)\)
      alone do not justify the filter when observation noises
      are correlated across time or have the wrong joint
      relationship with the hidden process.
\item If past realized costs reveal information that is excluded
      from the recorded observations, the displayed belief need
      not summarize the controller's actual information.
\item With an uninformative observation matrix, feedback can still
      make later beliefs depend on observations through the known
      chosen actions. The valid assertion is absence of new
      hidden-state information conditional on those actions.
\item An optimal policy for one initial law can behave
      suboptimally at unreachable beliefs.
\end{enumerate}
For finitely many unknown fixed transition parameters, augmenting
the hidden state by that parameter gives an exact finite model
and a posterior on the joint state and parameter. Replacing it
by a transition matrix averaged once at the initial prior
need not give the same predictions after data.
\end{hbtheorem}

\begin{proofboundary}{def:p17-astrom-source-boundary}{Finite horizons and later developments}
The 1965 paper proves finite-horizon results and gives two
four-stage examples. It supplies no infinite-horizon existence,
discounted convergence, average-reward, or approximation-by-grid
theorem. Its introductory discussion of recovering a continuous
model by limiting quantizations is a research direction, with
no error or convergence theorem established there.
The construction above optimizes on the entire belief simplex.
It has a finite observation branching rule but is not a finite
grid algorithm for that simplex.
The exact affine case calculations belong to the two specified
examples. \Cref{sec:p17-astrom-incomplete-information-ii} proves
the general finite-horizon convexity and concavity results of
the 1969 continuation, including controlled observation quality
and the homogeneous extension to the nonnegative cone.
General later results about finite piecewise-linear representations
retain their own source and proof obligations.
\end{proofboundary}
